% Drop-in replacement for Section V-D of
% "Formal Verification in Hybrid Quantum-Classical Computing Systems".
%
% Required preamble packages:
%   \usepackage{amsmath,amssymb,booktabs,tabularx,array,url}
% The paper already appears to load hyperref; if not, \usepackage{url} supplies
% \path for breakable repository paths.
%
% Integration warning (not intended for typesetting): the current numeric-layer
% subsection describes an older Estimator campaign (25 executed evaluations,
% 4,096 shots per term, 409,600 total shots).  This section reports the later
% primary SamplerV2 campaign (50 evaluations, four independent 1,024-shot PUBs
% per evaluation, 204,800 full-run shots).  The final manuscript should either
% label those as two distinct campaigns or update the older paragraph/table/
% figure so that readers do not interpret the numbers as contradictory reports
% of one experiment.

\subsection{Runtime Verification (DejaVu + Python Instrumentation)}
\label{sec:runtime-verification}

\subsubsection{Purpose and assurance boundary}

The runtime-verification layer checks whether a \emph{concrete, finite}
hybrid quantum--classical execution conforms to a declared execution protocol
while the execution is in progress, and whether the same verdict can be
reproduced later from durable evidence.  It is deliberately positioned at the
boundary between the classical controller, the quantum backend, the numeric
checker, and the optimizer.  It therefore observes facts that are abstracted
away by the NuXmv model and that cannot be established by a proof of an
isolated numeric kernel alone: the identity and order of parameter bindings,
backend-submission intent, returned job identifiers, term/PUB assignments,
raw-shot evidence, energy release, optimizer mutation, cumulative resource
use, and the declared terminal reason.

The principal claim supported by a passing run is conditional:

\begin{quote}
Given faithful instrumentation, correct cryptographic primitives, and the
archived observations, every observed event in the run obeyed the declared
binding, submission, measurement, numeric-certificate, optimizer-transition,
budget, convergence, and termination contracts, and independent replay
obtained the same verdict.
\end{quote}

This is a trace-conformance claim, not a proof that the cloud backend honestly
performed the claimed physical operations, that the device prepared the ideal
state, that the ansatz contains the true ground state, that the optimizer found
a global optimum, or that one stochastic campaign estimates hardware bias.
Likewise, a SHA-256 digest supplies integrity relative to a trusted digest and
archive; it is not a digital signature and does not by itself establish the
identity of the party that produced the data.

\subsubsection{Reference execution and observation model}

Let a run be identified by the pair $(r,c)$, where $r$ is a run identifier and
$c$ is the digest of the fully resolved configuration, including the bound
Hamiltonian manifest.  Let $k\in\{0,\ldots,K-1\}$ denote the zero-based
evaluation index, $e_k$ a fresh evaluation identifier, $\boldsymbol{\theta}_k$
the bound parameter vector, $s_k$ a submission identifier, $j_k$ the external
backend-job identifier, and $\widehat{E}_k$ the verified sampled objective.
The monitored execution is an ordered event sequence

\begin{equation}
  \tau = \langle a_1,a_2,\ldots,a_m\rangle,
  \qquad
  a_i=(i,t_i,n_i,\mathbf{x}_i,p_i),
  \label{eq:rv-trace}
\end{equation}

where $i$ is a strictly increasing sequence watermark, $t_i$ is a monotonic
host timestamp, $n_i$ is an event name, $\mathbf{x}_i$ is a tuple of
monitor-safe string atoms, and $p_i$ is a typed JSON payload used by the Python
contract engine and replay checker.  The compact DejaVu bridge transmits
\texttt{RV1}, the watermark, the event name, and the monitor arguments as a
tab-separated record.  Rich values such as floating-point energies, complete
parameter vectors, counts, and metadata remain in the JSON trace or linked raw
sidecar and are validated by Python.  Monitor atoms are restricted to the
regular language \texttt{[A-Za-z0-9.\_:-]+}; event names are similarly
restricted.  This prevents whitespace and delimiter ambiguity in the bridge
protocol.

The observed refinement of one nonterminal evaluation is

\begin{align}
  &\operatorname{start}(k,e_k,\boldsymbol{\theta}_k)
  \prec \operatorname{bind}(k,e_k,h^{\theta}_k)
  \prec \{\operatorname{termExpected}(k,e_k,s_k,P_i)\}_{i=0}^{3}
  \prec \operatorname{submitIntent}(k,e_k,s_k)
  \prec \operatorname{submitted}(k,e_k,s_k,j_k) \nonumber\\
  &\prec \{\operatorname{termObserved}(k,e_k,j_k,P_i)\}_{i=0}^{3}
  \prec \operatorname{measurementsSealed}(k,e_k)
  \prec \operatorname{numericValid}(k,e_k,h^{E}_k) \nonumber\\
  &\prec \operatorname{convergenceChecked}(k,e_k,h^{E}_k)
  \prec \operatorname{iterationPass}(k,e_k,h^{E}_k)
  \prec \operatorname{objectiveReleaseIntent}(k,e_k,h^{E}_k) \nonumber\\
  &\prec \operatorname{optimizerUpdateIntent}(k,e_k,h^{\theta}_{k+1})
  \prec \operatorname{optimizerUpdated}(k,e_k,h^{\theta}_{k+1})
  \prec \operatorname{iterationClosed}(k,e_k),
  \label{eq:rv-normal-order}
\end{align}

where $\prec$ denotes strict trace order and $h^{\theta}_k$ and $h^{E}_k$ are
canonical digests of the parameter vector and energy, respectively.  The two
optimizer events are omitted at a converged checkpoint or at the final
evaluation cap.  Every domain event is preceded by an \texttt{activity}
witness, allowing the run-level monitor to enforce that all activity follows
the matching start and that no activity occurs after a terminal event or after
\texttt{runClosed}.

The implementation is sequential: at most one evaluation is active, one
backend job is authorized for that evaluation, and the next evaluation cannot
start until the preceding evaluation is closed.  Concurrent evaluations,
batched optimizers with multiple outstanding proposals, multi-backend
execution, and distributed controllers are outside the current model.

\subsubsection{Synchronous, fail-closed enforcement}

The online gate is synchronous.  Before an externally visible side effect is
permitted, the controller (i) durably appends the proposed event, (ii) passes
the complete event to the Python contract engine, (iii) submits the compact
event to the generated DejaVu monitor, and (iv) requires a response carrying
the same sequence watermark.  If both checkers accept, the guarded operation
may proceed.  Formally, for a guarded operation $g_i$ associated with event
$a_i$,

\begin{equation}
  \operatorname{enabled}(g_i)
  \Longleftrightarrow
  \operatorname{persisted}(a_i)
  \land \operatorname{PythonOK}(a_i,\tau_{<i})
  \land \operatorname{DejaVuAck}(i,\mathrm{PASS}).
  \label{eq:rv-gate}
\end{equation}

A negative acknowledgement, watermark mismatch, malformed reply, monitor
timeout, monitor-process failure, trace-persistence failure, or Python contract
failure causes the path to fail closed.  The controller records
\texttt{criticalFail} when possible, emits \texttt{iterationAborted} if an
evaluation is active, emits \texttt{runAborted}, and does not release the
rejected value or perform the guarded mutation.  A missing or unavailable real
monitor produces an \emph{inconclusive} verdict; a test-only always-pass shadow
monitor is never allowed to mint temporal assurance.

The generated monitor is provenance-bound before startup.  The runtime checks
the SHA-256 digests of the QTL specification, Scala bridge, DejaVu runtime JAR,
and generated monitor JAR against a manifest embedded in the monitor artifact.
It also rejects a generated monitor that is stale relative to the checked-in
QTL or bridge source.  The reference run used a 24-bit finite domain per
monitor variable and a 30-second acknowledgement timeout.  The timeout is a
failure bound, not an expected latency; measured acknowledgement latencies are
reported below.

\subsubsection{Three levels of temporal contracts}

The QTL specification is written in past-time first-order temporal logic and
contains 68 generated-monitor properties.  Its audited classification contains
26 intra-iteration properties, 25 inter-iteration properties (including the
historically named \texttt{run\_iterationLimitOnce}), and 17 run-level
properties.  The Python checker separately registers 75 executable contract
identifiers: 25 intra-iteration, 27 inter-iteration, and 23 run-level.  Thus,
the explicit combined level registry contains 143 identifiers.  The concrete
hardware report records 37 dynamically exercised Python contract identifiers;
that run-dependent count must not be confused with the static property-set
size.

\begin{table*}[t]
  \centering
  \caption{Runtime-verification levels and representative obligations.}
  \label{tab:rv-levels}
  \renewcommand{\arraystretch}{1.12}
  \begin{tabular}{p{0.15\textwidth}p{0.22\textwidth}p{0.53\textwidth}}
    \toprule
    Level & Question & Representative obligations \\
    \midrule
    Intra-iteration
      & Was one evaluation executed consistently?
      & Parameters are bound after evaluation start and before submission;
        each declared term is observed exactly once and only after submission;
        sealing follows all required observations; numeric validity follows
        sealing; convergence is checked once; objective release follows an
        iteration pass; abort follows a recorded critical failure. \\
    Inter-iteration
      & Did the optimization history evolve coherently?
      & Evaluation, submission, raw-evidence, and backend-job identities are
        fresh; iterations are contiguous and nonoverlapping; the current
        convergence check names the last closed checkpoint; cumulative shots
        are continuous and bounded; optimizer intent follows objective release;
        commit follows intent exactly once; the successor binds the exact
        committed parameter digest. \\
    Run level
      & Did the finite run obey its lifecycle and declared terminal policy?
      & Exactly one matching start, terminal event, and close occur; all
        activity follows start; terminal states are absorbing; a normal terminal
        event has no prior critical failure; a failure forces abort; provenance,
        execution options, evidence completeness, and resource bounds match the
        configuration; the terminal reason is truthful. \\
    \bottomrule
  \end{tabular}
\end{table*}

Representative temporal implications are shown schematically below, where
$\mathbf{P}\phi$ means that the required witness occurs in the acknowledged
past and the strict-past form is used for duplicate exclusion:

\begin{align}
  \operatorname{jobSubmitIntent}(k,e,s,h^{\theta})
    &\Rightarrow \mathbf{P}\operatorname{parametersBound}(k,e,h^{\theta}),
    \label{eq:rv-qtl-submit}\\
  \operatorname{termObserved}(k,e,s,j,P)
    &\Rightarrow \mathbf{P}\operatorname{termExpected}(k,e,s,P)
      \land \mathbf{P}\operatorname{jobSubmitted}(k,e,s,j),
    \label{eq:rv-qtl-term}\\
  \operatorname{numericValid}(k,e,h^E)
    &\Rightarrow \mathbf{P}\operatorname{measurementsSealed}(k,e),
    \label{eq:rv-qtl-energy}\\
  \operatorname{optimizerUpdated}(k,e,h^{\theta}_{k+1})
    &\Rightarrow \mathbf{P}\operatorname{optimizerUpdateIntent}
      (k,e,h^{\theta}_{k+1}),
    \label{eq:rv-qtl-update}\\
  \operatorname{nextIterationStarted}(k+1,e_{k+1},e_k,h^{\theta}_{k+1})
    &\Rightarrow \mathbf{P}\operatorname{optimizerUpdated}
      (k,e_k,h^{\theta}_{k+1}),
    \label{eq:rv-qtl-next}\\
  \operatorname{runMaxIterations}(K-1,e_{K-1},h^E_{K-1})
    &\Rightarrow \mathbf{P}\operatorname{iterationLimitReached}(K-1,e_{K-1})
      \land \mathbf{P}\operatorname{convergenceFalse}(K-1,e_{K-1},h^E_{K-1})
      \land \mathbf{P}\operatorname{iterationClosed}(K-1,e_{K-1}).
    \label{eq:rv-qtl-terminal}
\end{align}

The Python and DejaVu components are intentionally complementary.  DejaVu is
the authority for causal order, strict uniqueness, and relational use of
identities over the observed history.  Python checks facts that are awkward or
unsafe to encode into the finite temporal domain: exact JSON schema and types,
finite-number requirements, vector dimensions and bounds, raw-shot decoding,
parity arithmetic, floating-point tolerances, Hamiltonian aggregation,
configuration immutability, cumulative resource arithmetic, and correlation
with the independently selected evidence sidecar.  Neither checker is treated
as a redundant implementation of the other.  Differential tests confirm this
division: some numeric/raw-data mutants are killed only by Python, while some
ordering and uniqueness mutants are killed only by DejaVu.

\subsubsection{Representative $\mathrm{H}_2$ SamplerV2 protocol}

The primary runtime-verification case study uses the same representative
two-qubit $\mathrm{H}_2$ model family as the other verification layers.  The
bound reduced Hamiltonian is

\begin{equation}
  H = \frac{1}{10^4}
  \left(3980\,YZ - 3980\,ZI - 113\,ZZ + 1810\,XX\right).
  \label{eq:rv-h2-hamiltonian}
\end{equation}

This is the objective of the archived four-term model.  It must not be silently
identified with a total molecular energy obtained using a different fermion
mapping, frozen-core convention, nuclear-repulsion constant, or identity-term
offset.  Because each Pauli expectation lies in $[-1,1]$, the manifest induces
the coarse but executable bound

\begin{equation}
  |\widehat{E}| \le
  \frac{3980+3980+113+1810}{10^4}=0.9883.
  \label{eq:rv-energy-bound}
\end{equation}

The ansatz is a two-qubit Qiskit \texttt{EfficientSU2} circuit with one
repetition, linear entanglement, \texttt{ry} single-qubit rotations, and four
parameters.  The initial parameter vector is generated by the declared seed
$7$:

\begin{equation}
  \boldsymbol{\theta}_0 =
  (-1.1068913836067251,\,-2.193779340387185,\,
    0.9483492633509023,\,-2.6864620414936398).
  \label{eq:rv-initial-theta}
\end{equation}

All parameters are required to be finite and within $[-\pi,\pi]$.  The
transpiler uses optimization level $3$ and seed $7$.  For the archived hardware
run, logical qubits $[0,1]$ were mapped to physical qubits $[6,7]$ on
\texttt{ibm\_marrakesh}; the same fixed mapping was recorded for every PUB.

Each objective evaluation consists of four \emph{independent} measured
circuits, submitted as four SamplerV2 PUBs in one IBM job.  Every PUB requests
and must return exactly $N=1024$ shots.  We intentionally do not group
$ZI$ and $ZZ$: consequently, ``1024 shots per term'' means 1024 physical
executions for each term rather than 1024 shared executions reported as two
effective term samples.  The physical cost of one evaluation and the declared
full-run ceiling are therefore

\begin{equation}
  N_{\mathrm{eval}}=4\times1024=4096,
  \qquad
  N_{\max}=50\times4\times1024=204{,}800.
  \label{eq:rv-shot-budget}
\end{equation}

\begin{table*}[t]
  \centering
  \caption{Independent measurement PUBs used in each $\mathrm{H}_2$ evaluation.  Qiskit count keys and Pauli labels are both ordered $q_1q_0$.}
  \label{tab:rv-measurements}
  \renewcommand{\arraystretch}{1.12}
  \begin{tabular}{ccp{0.34\textwidth}p{0.24\textwidth}r}
    \toprule
    PUB & Term & Basis change before $Z$ measurement & Negative outcomes & Shots \\
    \midrule
    0 & $YZ$ & $S^{\dagger}$ then $H$ on the $Y$ qubit; $Z$ unchanged
      & $01,10$ (odd parity) & 1024 \\
    1 & $ZI$ & No rotation; inspect $q_1$ and ignore $q_0$
      & $10,11$ & 1024 \\
    2 & $ZZ$ & No basis rotation
      & $01,10$ (odd parity) & 1024 \\
    3 & $XX$ & $H$ on both qubits
      & $01,10$ (odd parity) & 1024 \\
    \bottomrule
  \end{tabular}
\end{table*}

For a Pauli label $P=P_{n-1}\cdots P_0$ and a returned bitstring
$b=b_{n-1}\cdots b_0$, define the active-qubit set

\begin{equation}
  A(P)=\{q\mid P_q\neq I\},
  \qquad
  \lambda_P(b)=(-1)^{\sum_{q\in A(P)}b_q}.
  \label{eq:rv-parity-sign}
\end{equation}

The complete count dictionary $C_P(b)$ is converted into a negative-eigenvalue
count and expectation by

\begin{align}
  N_{-}(P) &= \sum_b C_P(b)\,
    \mathbf{1}\!\left[\lambda_P(b)=-1\right],
    \label{eq:rv-negative-count}\\
  \widehat{\mu}_P &= \frac{N_{+}(P)-N_{-}(P)}{N}
    = 1-\frac{2N_{-}(P)}{N},
    \qquad N=1024.
    \label{eq:rv-expectation}
\end{align}

Here \texttt{counts\_ones} is treated as the explicit alias of
$N_{-}(P)$, not as the number of one-bits on an arbitrarily chosen qubit.  The
sampled objective is reconstructed using the integer-scaled coefficients:

\begin{equation}
  \widehat{E}_k = \frac{1}{10^4}
  \left(3980\widehat{\mu}_{YZ,k}
  -3980\widehat{\mu}_{ZI,k}
  -113\widehat{\mu}_{ZZ,k}
  +1810\widehat{\mu}_{XX,k}\right).
  \label{eq:rv-energy-reconstruction}
\end{equation}

The Python gate requires the declared per-term expectation to equal
Eq.~\eqref{eq:rv-expectation} and the declared energy to equal
Eq.~\eqref{eq:rv-energy-reconstruction} with zero relative tolerance and
absolute tolerance $10^{-12}$.  The raw-evidence loader uses a stricter
$10^{-15}$ absolute tolerance when comparing the sidecar's expectation with
the value directly reconstructed from its integer parity count.

\subsubsection{Raw evidence, provenance, and independent replay}

For each PUB, the runtime extracts the named Qiskit \texttt{BitArray} and
archives (i) the requested and observed shot totals, (ii) all ordered
bitstrings, (iii) the complete count dictionary, (iv) the exact C-order
\texttt{uint8} packed array encoded as base64, (v) the Pauli label and bit-order
convention, (vi) $N_-$ and $\widehat{\mu}_P$, (vii) evaluation, term, PUB,
submission, and IBM-job identities, (viii) bound parameters and their digest,
(ix) parameterized and executable-circuit QPY digests, (x) physical layout,
transpiler seed and optimization level, (xi) timestamps, and (xii) returned
job/PUB metadata.  The packed representation must decode byte-for-byte to the
ordered bitstrings; the ordered bitstrings must regenerate the count
dictionary; and the count values must sum to the observed shot total.

Let $R_i$ be a canonical JSON PUB record with its own digest field removed.
Its content digest is

\begin{equation}
  d_i=\operatorname{SHA256}(\operatorname{CanonicalJSON}(R_i)).
  \label{eq:rv-record-digest}
\end{equation}

The sidecar is appended and flushed (with \texttt{fsync} in durable mode)
before the compact \texttt{termObserved} event carrying $d_i$ is emitted.  At
the end of the run, a manifest binds the run identifier, configuration digest,
record count, cumulative shots, ordered list $(d_0,\ldots,d_{199})$, and an
aggregate digest

\begin{equation}
  d_{\mathrm{records}} =
  \operatorname{SHA256}(\operatorname{CanonicalJSON}([d_0,\ldots,d_{199}])).
  \label{eq:rv-manifest-aggregate}
\end{equation}

The manifest is itself canonically hashed with its own digest field omitted.
Duplicate hashes, missing indices, records after the manifest, mismatched run
or configuration identities, incorrect cumulative totals, or an altered
record order are rejected.  The compact trace is a strict-schema, ordered
JSONL event journal and is separately archived with a whole-file SHA-256; it is
not claimed to be a cryptographic hash chain.

Replay takes the trace, configuration, and evidence sidecar as three explicit
inputs.  In particular, replay does not follow an evidence path named by the
trace under test.  This prevents a manipulated trace from selecting its own
replacement evidence archive.  Replay reconstructs every packed shot array,
ordered-shot list, count dictionary, parity, expectation, energy, cumulative
budget, ask/tell transition, and terminal certificate, while re-executing the
same temporal event stream through a provenance-checked DejaVu monitor.  A
truncated trace, a missing sidecar, a shadow monitor, or an unavailable monitor
cannot yield \textsc{pass}; the result is \textsc{inconclusive} unless a
specific violated contract yields \textsc{fail}.

\subsubsection{Verifier-mediated ask/tell optimization}

An ordinary callback-style optimizer can choose or mutate its internal state
before a monitor learns what happened.  That is observational monitoring, not
pre-authorization.  The primary experiment therefore uses an explicit
ask/tell pattern-search optimizer whose preview operation is pure.  Let $z_k$
be the private optimizer state after checkpoint $k$.  Preview satisfies

\begin{equation}
  (\widetilde{\boldsymbol{\theta}}_{k+1},u_k)
    =\operatorname{Preview}(z_k,\boldsymbol{\theta}_k,\widehat{E}_k),
  \qquad z_k\ \text{unchanged},
  \label{eq:rv-preview}
\end{equation}

where $u_k$ is a single-use opaque token bound to the proposal, generation,
and candidate vector.  The verifier then emits and synchronously acknowledges

\begin{equation}
  \operatorname{optimizerUpdateIntent}
  (r,c,k,e_k,h(\widetilde{\boldsymbol{\theta}}_{k+1})).
  \label{eq:rv-authorize}
\end{equation}

Only after that acknowledgement may the controller call

\begin{equation}
  z_{k+1}=\operatorname{Commit}
  (z_k,u_k,\texttt{authorized\_theta}=\widetilde{\boldsymbol{\theta}}_{k+1}).
  \label{eq:rv-commit}
\end{equation}

Commit rejects a foreign, altered, stale, reused, or wrong-generation token,
and rejects an authorized vector different from the previewed vector.  After a
successful mutation, \texttt{optimizerUpdated} records the same digest.  The
next evaluation is accepted only if it binds that exact digest:

\begin{equation}
  h(\boldsymbol{\theta}_{k+1})
  =h(\widetilde{\boldsymbol{\theta}}_{k+1})
  =h^{\theta}_{k+1}.
  \label{eq:rv-theta-continuity}
\end{equation}

The pattern search uses initial step $0.4$, step shrink factor $0.5$, zero
improvement tolerance, coordinate-wise candidates clipped to $[-\pi,\pi]$,
and deterministic initialization from seed $7$.  This runtime contract does
not prove that pattern search is an optimal VQE optimizer.  It proves the
narrower but operationally essential property that the verifier sees and
accepts the exact candidate before the only state mutation, and that the next
quantum evaluation uses the committed candidate.

For evaluation $0$, the archived trace records

\begin{align}
  \widehat{E}_0 &= 0.1916484375,\nonumber\\
  \widetilde{\boldsymbol{\theta}}_1
    &=(-0.7068913836,-2.1937793404,0.9483492634,-2.6864620415),
  \label{eq:rv-first-proposal}
\end{align}

up to the displayed precision.  The proposal was authorized, committed once,
and bound by evaluation $1$ with the same full-precision digest.  Across the
50-evaluation hardware trace, replay found 49 proposal intents, 49 commits,
and 49 exact successor bindings in strict order.

\subsubsection{Convergence and terminal semantics}

The runtime, rather than the optimizer, computes the declared convergence
predicate.  Let $\epsilon=0.005$, $k_{\min}=10$, and $p=3$.  For $k>0$ define

\begin{align}
  \Delta_k &= |\widehat{E}_k-\widehat{E}_{k-1}|,
  \label{eq:rv-delta}\\
  s_k &=
  \begin{cases}
    s_{k-1}+1, & \Delta_k\le \epsilon+\epsilon_E,\\
    0, & \text{otherwise},
  \end{cases}
  \qquad s_0=0,
  \label{eq:rv-stable-count}\\
  C_k &=
  \left(\Delta_k\le\epsilon+\epsilon_E\right)
  \land(k+1\ge k_{\min})\land(s_k\ge p),
  \label{eq:rv-convergence}
\end{align}

where $\epsilon_E=10^{-12}$ is the configured absolute numeric tolerance.
The first evaluation is never declared converged.  The event
\texttt{nextConvergenceChecked} carries the current evaluation, energy and
parameter identities together with the immediately preceding evaluation and
energy identity, and Python recomputes $\Delta_k$, $s_k$, and $C_k$.

Exactly one of the following terminal reasons is permitted:

\begin{enumerate}
  \item \texttt{converged}, only after a closed checkpoint with $C_k=\mathrm{true}$;
  \item \texttt{max\_iterations}, only after evaluation $K-1$ is closed with
        $C_{K-1}=\mathrm{false}$ and a unique iteration-limit certificate;
  \item \texttt{optimizer\_stopped}, only when explicitly allowed, before the
        cap, after a closed nonconverged checkpoint; or
  \item \texttt{aborted}, only after a recorded critical failure.
\end{enumerate}

The H$_2$ Sampler profiles set $K=50$ and disallow optimizer-stop termination.
The reference hardware and statevector runs both exhausted the cap without
satisfying Eq.~\eqref{eq:rv-convergence}; therefore
\texttt{max\_iterations}, rather than \texttt{converged}, is the only truthful
normal terminal reason.

\subsubsection{Fixed execution policy}

\begin{table*}[t]
  \centering
  \caption{Declared and observed configuration of the primary hardware campaign.}
  \label{tab:rv-fixed-config}
  \renewcommand{\arraystretch}{1.1}
  \begin{tabular}{p{0.24\textwidth}p{0.31\textwidth}p{0.35\textwidth}}
    \toprule
    Field & Declared value & Observed/reference value \\
    \midrule
    System and primitive
      & Two-qubit $\mathrm{H}_2$; IBM Runtime SamplerV2
      & Same reduced four-term model and primitive \\
    Backend
      & \texttt{ibm\_marrakesh}, subject to availability
      & \texttt{ibm\_marrakesh}; backend version 1.0.21 at pre-audit; 156 qubits \\
    Evaluations and PUBs
      & At most 50; four independent PUBs per evaluation
      & 50 evaluations; 200 PUBs; 50 unique IBM jobs \\
    Shots
      & 1024 per PUB; 4096 per evaluation; 204,800 full-run maximum
      & Exactly 1024 for every PUB; 204,800 full-run shots \\
    Term/PUB order
      & $YZ,ZI,ZZ,XX$ at indices $0,1,2,3$
      & Exact match for all evaluations \\
    Optimizer
      & Ask/tell pattern search; authorization before commit
      & 49 intent/commit/successor-binding chains verified \\
    Convergence
      & $\epsilon=0.005$, minimum 10 evaluations, patience 3
      & Not reached; cap certificate accepted \\
    Layout
      & One transpiled layout fixed across all four circuit families and all evaluations
      & Logical $[0,1]$ to physical $[6,7]$ \\
    Numeric tolerances
      & Energy and expectation absolute tolerance $10^{-12}$
      & All reconstructions within tolerance \\
    \bottomrule
  \end{tabular}
\end{table*}

The Sampler options are explicitly recorded rather than inherited from future
defaults: \texttt{default\_shots=1024}, run-level \texttt{shots=1024},
classified measurement output, gate twirling disabled, measurement twirling
disabled, and dynamical decoupling disabled.  Measurement mitigation, zero-noise
extrapolation (ZNE), and probabilistic error cancellation (PEC) are marked not
applicable to this direct raw-count SamplerV2 path.  These choices preserve the
auditable transformation

\begin{equation}
  \text{IBM bitstrings}\to\text{counts}\to N_-
  \to\widehat{\mu}_P\to\widehat{E}_k
  \to\text{verified release}\to\text{optimizer}.
  \label{eq:rv-evidence-chain}
\end{equation}

They do not mitigate device or readout error.  The hardware objective should
therefore be interpreted as a raw, auditable sampled objective, not as a
chemically corrected energy estimate.

\subsubsection{Validation before hardware execution}

Before the primary hardware submission, a local validator executed 36 checks
with zero failures.  It included a one-qubit Hadamard circuit with exactly 1024
outcomes, all four H$_2$ basis-measurement circuits through
\texttt{StatevectorSampler}, comparison of count-derived expectations with
\texttt{StatevectorEstimator}, explicit confirmation of Qiskit's
$q_{n-1}\cdots q_0$ Pauli-label/count-key convention, and negative tests for a
missing PUB, wrong term/PUB assignment, shot mismatch, corrupted counts,
corrupted ordered shots, and corrupted packed data.  The pre-submission
automated suite passed 129 tests in 14.787 seconds; a final post-artifact rerun
again passed all 129 tests in 14.322 seconds.

The one-evaluation IBM preflight used job
\texttt{dafnocu42tqs73b06mb0} and 4096 physical shots.  Hardware-specific
result fields, counts, parity, and layout checks succeeded, but independent
replay rejected a circuit-provenance digest.  Investigation found that the same
parameterized circuit had been serialized to QPY twice and the two byte streams
were not stable across serializations.  The correction captures each QPY
serialization once and reuses its digest for both evidence and run provenance;
a regression test enforces that rule.  The failed preflight artifact was
preserved and was not resubmitted.  The complete 50-evaluation run used the
corrected implementation and passed replay.  This incident is evidence that
the provenance contract is active: it rejected inconsistent identity evidence
even though the raw hardware counts were otherwise valid.

\subsubsection{Primary hardware result}

The reference campaign, run identifier
\texttt{h2-ibm-marrakesh-sampler-20260907}, completed on
\texttt{ibm\_marrakesh} on September 7--8, 2026.  It produced 2103 compact
trace events, 200 full PUB sidecar records, and one sealed evidence manifest.
Online verification and independent replay both returned \textsc{pass} at the
intra-iteration, inter-iteration, and run levels with zero findings.

\begin{table*}[t]
  \centering
  \caption{Observed results of the primary hardware campaign.}
  \label{tab:rv-hardware-results}
  \renewcommand{\arraystretch}{1.1}
  \begin{tabular}{p{0.27\textwidth}p{0.18\textwidth}p{0.27\textwidth}p{0.18\textwidth}}
    \toprule
    Metric & Result & Metric & Result \\
    \midrule
    Evaluations & 50 & Unique IBM jobs & 50 \\
    Independent PUB records & 200 & Physical shots & 204,800 \\
    Compact trace events & 2103 & Online/replay findings & 0 / 0 \\
    Online verdict & \textsc{pass} & Replay verdict & \textsc{pass} \\
    Intra-iteration level & \textsc{pass} & Inter-iteration level & \textsc{pass} \\
    Run level & \textsc{pass} & Terminal reason & \texttt{max\_iterations} \\
    Initial objective & $0.1916484375$ & Final objective & $-0.3775728515625$ \\
    Best objective & $-0.418932421875$ & Best evaluation & 28 (zero-based) \\
    Monitored wall duration & $1918.285$ s & IBM QPU charge time & $150$ s \\
    Estimated quantum time & $137.703$ s & Physical layout & $[6,7]$ \\
    \bottomrule
  \end{tabular}
\end{table*}

At the best hardware checkpoint, evaluation 28 and IBM job
\texttt{dafoahbdd5gc73da3qf0}, the four count dictionaries reconstruct as shown
in Table~\ref{tab:rv-best-checkpoint}.

\begin{table*}[t]
  \centering
  \caption{Count-level reconstruction of the best observed hardware checkpoint.  Count columns are ordered $00/01/10/11$.}
  \label{tab:rv-best-checkpoint}
  \renewcommand{\arraystretch}{1.1}
  \begin{tabular}{clrrr}
    \toprule
    Term & Counts $00/01/10/11$ & $N_-$ & $\widehat{\mu}_P$ & Weighted contribution \\
    \midrule
    $YZ$ & $158/387/162/317$ & 549 & $-0.072265625$ & $-0.028761718750$ \\
    $ZI$ & $350/660/5/9$       & 14  & $ 0.972656250$ & $-0.387117187500$ \\
    $ZZ$ & $314/688/6/16$      & 694 & $-0.355468750$ & $ 0.004016796875$ \\
    $XX$ & $479/32/500/13$     & 532 & $-0.039062500$ & $-0.007070312500$ \\
    \midrule
    \multicolumn{4}{r}{Reconstructed objective} & $-0.418932421875$ \\
    \bottomrule
  \end{tabular}
\end{table*}

For example, $N_-(YZ)=C(01)+C(10)=387+162=549$, so
$\widehat{\mu}_{YZ}=1-2(549)/1024=-0.072265625$.  Substitution into
Eq.~\eqref{eq:rv-energy-reconstruction} gives

\begin{align}
  \widehat{E}_{28}
  &=0.3980(-0.072265625)-0.3980(0.972656250)\nonumber\\
  &\quad-0.0113(-0.355468750)+0.1810(-0.039062500)\nonumber\\
  &=-0.418932421875,
  \label{eq:rv-best-energy}
\end{align}

which matches the trace value within the declared tolerance.  The offline
analyzer repeated the same reconstruction for all 200 PUBs and all 50 energy
checkpoints.

The full-run 204,800-shot ceiling excludes the separate 4096-shot preflight.
The total physical-shot consumption of the archived campaign including that
preflight was therefore

\begin{equation}
  204{,}800+4{,}096=208{,}896.
  \label{eq:rv-campaign-shots}
\end{equation}

No inference about adherence to the full-run ceiling should incorrectly count
the explicitly separate preflight against the primary-run configuration.

\subsubsection{Comparable shot-based statevector result}

The statevector comparator uses the same ansatz, initial vector, four
measurement circuits, term/PUB ordering, 1024 shots per term, 50-evaluation
cap, ask/tell optimizer, convergence rule, raw-evidence schema, and online and
replay monitors.  It uses \texttt{StatevectorSampler} to generate sampled
bitstrings; it is therefore not a shot-free exact objective.  This makes the
measurement and budget semantics directly comparable with hardware.

\begin{table*}[t]
  \centering
  \caption{Descriptive comparison of the two sampled optimization paths.}
  \label{tab:rv-hw-sv-comparison}
  \renewcommand{\arraystretch}{1.1}
  \begin{tabular}{p{0.34\textwidth}rr}
    \toprule
    Measure & IBM hardware & Statevector sampler \\
    \midrule
    Evaluations / PUBs / samples & $50/200/204{,}800$ & $50/200/204{,}800$ \\
    Initial objective & $0.1916484375$ & $0.2069205078125$ \\
    Best observed objective & $-0.418932421875$ & $-0.439887500000$ \\
    Best evaluation (zero-based) & 28 & 40 \\
    Final objective & $-0.3775728515625$ & $-0.439807421875$ \\
    Online / replay verdict & \textsc{pass} / \textsc{pass} & \textsc{pass} / \textsc{pass} \\
    Terminal reason & \texttt{max\_iterations} & \texttt{max\_iterations} \\
    \bottomrule
  \end{tabular}
\end{table*}

The difference between the best observed objectives is

\begin{equation}
  \widehat{E}^{\star}_{\mathrm{HW}}
  -\widehat{E}^{\star}_{\mathrm{SV}}
  =0.020955078125.
  \label{eq:rv-best-difference}
\end{equation}

This is a descriptive result for two adaptive sampled trajectories, not an
estimate of device bias at a fixed parameter vector.  The first difference in
sampled energy can change the next ask/tell decision, after which the two runs
may evaluate different parameter vectors.  A hardware-accuracy study would
instead require a preregistered set of fixed vectors, repeated independent
hardware campaigns or paired jobs, calibration covariates, uncertainty
intervals, and a declared statistical analysis plan.

\subsubsection{Acceptance criteria and verdict aggregation}

The hardware result is accepted only when every condition in
Table~\ref{tab:rv-acceptance} holds.  In the archived campaign all conditions
evaluated to true.

\begin{table*}[t]
  \centering
  \caption{Acceptance criteria for the primary runtime-verification result.}
  \label{tab:rv-acceptance}
  \renewcommand{\arraystretch}{1.1}
  \begin{tabular}{p{0.53\textwidth}p{0.29\textwidth}c}
    \toprule
    Requirement & Evidence/check & Result \\
    \midrule
    Every term has exactly 1024 observed outcomes
      & All 200 sidecar records and \texttt{BitArray.num\_shots} & \textsc{pass} \\
    Every complete count dictionary sums to 1024
      & Independent reconstruction from ordered shots & \textsc{pass} \\
    Packed bytes, ordered shots, counts, and hashes agree
      & Strict sidecar loader and sealed manifest & \textsc{pass} \\
    PUB indices and term identities are complete, unique, and ordered
      & Expected manifest versus all observed PUBs & \textsc{pass} \\
    Every expectation reconstructs from odd active-qubit parity
      & Eq.~\eqref{eq:rv-expectation}, all 200 PUBs & \textsc{pass} \\
    Every energy reconstructs from four expectations
      & Eq.~\eqref{eq:rv-energy-reconstruction}, all 50 checkpoints & \textsc{pass} \\
    Full-run cumulative shots never exceed 204,800
      & Per-checkpoint continuity and manifest total & \textsc{pass} \\
    Online and independent replay verdicts agree
      & Three levels in both JSON reports & \textsc{pass} \\
    Ask/tell authorization precedes each optimizer mutation
      & 49 ordered intent/commit/successor chains & \textsc{pass} \\
    Backend job and evaluation identities are fresh
      & 50 unique IBM job IDs and 50 evaluation IDs & \textsc{pass} \\
    Physical layout is fixed across executable circuits
      & All hardware records report logical $[0,1]\mapsto[6,7]$ & \textsc{pass} \\
    Evidence manifest commits to all records and the exact total
      & 200 ordered hashes; 204,800 observed shots & \textsc{pass} \\
    Terminal reason agrees with the final verified checkpoint
      & Nonconverged evaluation 49 plus unique cap certificate & \textsc{pass} \\
    \bottomrule
  \end{tabular}
\end{table*}

At each level, a recorded violation yields \textsc{fail}; missing assurance
that prevents the property from being established yields
\textsc{inconclusive}.  A level is \textsc{pass} only after the trace is
complete and no property assigned to that level is failed or inconclusive.
The overall report is \textsc{pass} only if all three levels pass and the run
is complete:

\begin{equation}
  V(\tau)=\mathrm{PASS}
  \Longleftrightarrow
  \operatorname{complete}(\tau)
  \land\bigwedge_{\ell\in\{\mathrm{intra,inter,run}\}}
  V_{\ell}(\tau)=\mathrm{PASS}.
  \label{eq:rv-verdict}
\end{equation}

This prevents an empty, truncated, monitorless, or only partially checked trace
from being misreported as successful merely because no explicit violation was
encountered.

\subsubsection{Executable overhead and scaling evaluation}

Runtime overhead is measured independently of hardware queue and execution
time.  The executable harness feeds identical deterministic H$_2$ observations
to (i) an uninstrumented in-memory loop, (ii) the Python runtime session with a
test/shadow monitor, and (iii) the Python runtime session with the real
synchronous generated DejaVu monitor.  The primary comparison uses five
warm-ups and 30 measured repetitions per condition.  It archives every run,
every acknowledgement latency, and periodic Python/JVM resident-set-size (RSS)
sample in CSV; JSON contains derived summaries only.

For the three-evaluation primary comparison, the generated monitor had
$358.260$ ms median and $375.798$ ms p95 end-to-end time, including a
$228.539$ ms median JVM startup.  Across 3810 acknowledged events, gate latency
was $0.505$ ms median and $1.722$ ms p95.  The corresponding Python-shadow p95
was $7.505$ ms end-to-end, and the intentionally minimal uninstrumented loop
had $0.050$ ms p95.  Relative percentages against that tiny baseline are
dominated by fixed JVM startup and must not be generalized to a hardware VQE.

The scaling sweep fixes the H$_2$ event protocol while varying the number of
nonconverging evaluations.  Table~\ref{tab:rv-scaling} reports the raw-summary
values from one measured scaling run at each size.

\begin{table*}[t]
  \centering
  \caption{Generated-monitor scaling with evaluation-history length.}
  \label{tab:rv-scaling}
  \renewcommand{\arraystretch}{1.08}
  \begin{tabular}{rrrrrr}
    \toprule
    Evaluations & Events & \shortstack{End-to-end\\(s)}
      & \shortstack{All-gate\\p95 (ms)}
      & \shortstack{\texttt{nextConvergenceChecked}\\p95 (ms)}
      & \shortstack{JVM peak RSS\\(MiB)} \\
    \midrule
     10 &  423 &  0.537 &  1.214 &  1.332 & 219.03 \\
     25 & 1053 &  0.819 &  1.063 &  1.237 & 273.66 \\
     50 & 2103 &  1.210 &  0.984 &  1.305 & 439.80 \\
     60 & 2523 &  1.350 &  0.929 &  1.669 & 481.05 \\
    100 & 4203 &  2.081 &  0.970 &  1.532 & 635.36 \\
    125 & 5253 &  2.929 &  1.438 &  2.307 & 718.81 \\
    150 & 6303 &  6.436 &  3.722 &  4.373 & 839.92 \\
    175 & 7353 & 12.486 &  7.742 &  7.606 & 847.94 \\
    200 & 8403 & 29.971 & 16.362 & 16.222 & 941.80 \\
    \bottomrule
  \end{tabular}
\end{table*}

The monitor is not constant-cost over unbounded history: resident memory grows
throughout the sweep, and a pronounced latency increase appears beyond roughly
125 evaluations.  At the fixed primary bound of 50, however, all-gate p95 is
$0.984$ ms and \texttt{nextConvergenceChecked} p95 is $1.305$ ms, far below
the 30-second hardware acknowledgement timeout.

\paragraph{Decision on \texttt{nextConvergenceChecked}.}
The current relation is retained for the reported 50-evaluation profile.  It
binds the current evaluation, current energy and parameter identities, previous
evaluation, previous energy identity, and convergence verdict.  Projecting
fields out merely to reduce monitor state would weaken the predecessor and
energy-continuity guarantee without materially improving the bounded campaign.
The defensible conclusion is therefore bounded: the relation is acceptable at
$K\le 50$ under the measured environment; these data do not justify unbounded
scaling.

Before materially increasing the profile beyond 100 evaluations, the planned
redesign is to let Python issue a compact checkpoint certificate, for example

\begin{equation}
  \chi_k=\operatorname{SHA256}\!\left(
  c\parallel e_k\parallel h^{\theta}_k\parallel h^E_k
  \parallel e_{k-1}\parallel h^E_{k-1}\parallel s_k\right),
  \label{eq:rv-future-certificate}
\end{equation}

after verifying the exact energy delta, patience state, theta continuity, and
predecessor relation.  DejaVu would then relate only the compact certificate
identity and causal order.  This proposal is not part of the reported system.
It must be accepted only after (i) differential mutants demonstrate unchanged
rejection of stale predecessor, stale energy, false convergence, and reordered
check faults, (ii) the generated monitor is rebuilt from pinned dependencies,
and (iii) the warm-up/repetition benchmark and raw RSS/latency sweep are rerun
through at least 400 evaluations.

\subsubsection{Fault model and negative testing}

The implementation is tested against faults at each assurance boundary.  The
negative suite includes, among others:

\begin{itemize}
  \item binding outside the parameter domain, wrong parameter dimension,
        NaN/Inf values, stale theta, and a successor that differs from the
        authorized optimizer proposal;
  \item term observation before declaration or submission, a missing PUB,
        duplicate/wrong PUB index, wrong Pauli assignment, mismatched
        negative-parity alias, and incomplete measurement sealing;
  \item requested/observed shot mismatch, counts not summing to the shot total,
        counts differing from ordered bitstrings, packed bytes differing from
        ordered bitstrings, and attacker-rehashed raw-data corruption;
  \item incorrect expectation, incorrect weighted contribution, incorrect
        total energy, out-of-bound energy, reused raw-evidence digest, and an
        evidence manifest with missing, duplicated, reordered, or orphaned
        records;
  \item duplicate evaluation, submission, or backend-job identity; skipped or
        overlapping evaluations; a convergence check using a non-last
        checkpoint; cumulative-budget discontinuity; and configuration drift;
  \item optimizer mutation before objective release, commit without matching
        intent, token alteration/reuse, missing commit, double commit, update at
        the iteration cap, and a next binding that differs from the commit;
  \item false convergence, convergence followed by a successor, iteration-cap
        termination without a cap certificate, optimizer-stop disguised as a
        cap, activity after terminal/close, multiple terminals, and normal
        completion after a critical failure;
  \item unavailable/stale monitor artifacts, monitor watermark mismatch,
        monitor timeout, real DejaVu rejection on the online path, empty or
        truncated trace, and replay without an explicitly supplied sidecar.
\end{itemize}

The tests establish fault-detection behavior for the modeled mutants; they do
not establish a complete fault model for all Qiskit, IBM Runtime, operating
system, network, or hardware failures.

\subsubsection{Assumptions, trust boundaries, and threats to validity}

The runtime claim depends on the following explicit assumptions.

\paragraph{Instrumentation completeness.}
All relevant controller-side effects must pass through the instrumented
session.  Direct calls to the backend or optimizer outside the gate are not
observable and therefore not constrained.  The Python process, event-journal
implementation, and bridge are included in the trusted computing base.

\paragraph{Sequential semantics.}
The formal event model assumes one run, one active evaluation, one authorized
job per evaluation, and one committed successor.  It does not cover concurrent
job submission, speculative optimizers, asynchronous cancellation races, or
recovery across multiple controller processes.

\paragraph{Backend cooperation.}
IBM job identifiers, timestamps, result metadata, and returned bitstrings are
treated as observations supplied by a cooperative service.  The verifier
checks internal consistency and provenance of what was returned; it does not
attest the physical device, prevent a malicious backend from fabricating a
self-consistent result, or verify IBM's accounting implementation.

\paragraph{Cryptographic and archival assumptions.}
SHA-256 is assumed collision resistant.  Canonical JSON, stored file digests,
and the archive channel are assumed to be implemented and retained correctly.
Because the evidence is hashed rather than digitally signed, an adversary able
to replace both an artifact and every trusted external digest remains outside
the threat model.

\paragraph{Software semantics and versioning.}
The result is tied to CPython 3.12.14, Qiskit 2.5.2,
\texttt{qiskit-ibm-runtime} 0.49.0, NumPy 2.5.3, SciPy 1.18.1, and OpenJDK
17.0.4.1 on arm64 macOS 26.6.2.  A new version may change QPY serialization,
bit-array structure, transpilation, default options, floating-point behavior,
or optimizer order; the validation, provenance, and replay suites must be
rerun rather than assuming compatibility.

\paragraph{Floating-point boundary.}
Counts and Hamiltonian coefficient integers are exact, but expectation and
energy values cross a binary floating-point interface.  The runtime uses
explicit finite-number checks, zero relative tolerance, integer-scaled
coefficients, and declared absolute tolerances.  It does not formally verify
the complete floating-point implementation of Qiskit, Python, NumPy, or the
optimizer.

\paragraph{Measurement interpretation.}
The Pauli label and Qiskit count key are assumed to share the documented
$q_{n-1}\cdots q_0$ order.  All qubits are measured; identity positions are
retained in the raw record but excluded from the parity calculation.  The
preflight and statevector validation test this convention for the four
reported terms, but not for arbitrary registers, dynamic circuits, mid-circuit
measurements, or parameter-swept PUB shapes.

\paragraph{Statistical scope.}
The run is one adaptive stochastic trajectory.  Passing trace contracts says
nothing about chemical accuracy, estimator variance, optimizer success
probability, cross-day calibration drift, or generalization to repeated
campaigns, other backends, deeper ansatze, larger Hamiltonians, or error
mitigation.  Those questions require the probabilistic layer and a separate
experimental design.

\paragraph{Construct validity of the benchmark.}
The overhead harness uses deterministic synthetic observations to isolate
verification costs.  It intentionally excludes network, queue, and device
time.  Its uninstrumented baseline is extremely small, so percentage overhead
is not representative of a real VQE wall-clock fraction.  Absolute gate
latency and RSS under the stated environment are the appropriate reported
quantities.  The one-run-per-scale sweep reveals a clear late-history trend but
does not provide confidence intervals for every scale.

\paragraph{Physical layout and calibration.}
The archived run used physical qubits $[6,7]$, but backend topology,
calibration, availability, queue conditions, and selected layout can change.
A reproduction may use a different valid fixed layout and will generally
produce different counts.  Reproduction means matching the declared protocol,
evidence completeness, and verdict semantics, not reproducing exact random
bitstrings or the energy sequence.

\subsubsection{Reproducibility and durable artifacts}

The archived environment manifest records package versions, commands, source
and artifact SHA-256 values, backend audit information, and the selected
layout.  The principal machine-readable artifacts are:

\begin{itemize}
  \item configuration: \path{configs/h2_ibm_sampler_50.toml};
  \item comparable statevector configuration:
        \path{configs/h2_statevector.toml};
  \item Hamiltonian manifest: \path{data/h2_hamiltonian.json};
  \item temporal specification: \path{specs/vqe_runtime.qtl};
  \item hardware trace:
        \path{artifacts/traces/h2_ibm_marrakesh_sampler_20260907.jsonl};
  \item hardware raw evidence:
        \path{artifacts/evidence/h2_ibm_marrakesh_sampler_20260907.jsonl};
  \item online and replay reports:
        \path{artifacts/reports/h2_ibm_marrakesh_sampler_20260907_online.json}
        and
        \path{artifacts/reports/h2_ibm_marrakesh_sampler_20260907_replay.json};
  \item statevector trace/evidence/reports with the corresponding
        \path{h2_statevector_sampler_20260907} basename;
  \item comparison summary and raw per-evaluation/per-PUB CSVs:
        \path{artifacts/results/h2_sampler_comparison_20260907/};
  \item executable benchmark and raw benchmark output:
        \path{tools/benchmark_overhead.py} and
        \path{artifacts/benchmarks/}; and
  \item exact environment and digest manifest:
        \path{artifacts/environments/h2_sampler_campaign_20260907.json}.
\end{itemize}

For the hardware run, the trace SHA-256 is
\texttt{c954babd7d629f8e\allowbreak{}75f5cf5f69132108\allowbreak{}
aa99c6cc6bb04a18\allowbreak{}fb046fb0faeb342c},
the evidence-file SHA-256 is
\texttt{e93da9837ca8c504\allowbreak{}f8d3120ddf97d6bc\allowbreak{}
048f923f56f6ed77\allowbreak{}eaa3a12509e50671},
and the sealed evidence-manifest SHA-256 is
\texttt{06019a3a2fff0838\allowbreak{}8e13a0ebce7faa1f\allowbreak{}
e2f10203ebc83049\allowbreak{}7bbc1bc9078fdd57}.
The online-report and replay-report SHA-256 values are, respectively,
\texttt{4ef8d86760c2a81c\allowbreak{}fd1a5a0c2a0e360f\allowbreak{}
9d8284a68f09569c\allowbreak{}4ac43678ecb6f4cf}
and
\texttt{a94c0613f6b151db\allowbreak{}e4fb0031197b4d69\allowbreak{}
8f7ea7b61bf305a1\allowbreak{}5be35b9649ac1f51}.

In summary, the runtime layer turns the abstract ordering requirement
``no optimizer update before verified energy aggregation'' into an enforced
end-to-end protocol.  Raw IBM bitstrings are durably linked to parity-derived
expectations and a reconstructed H$_2$ objective; that objective is released
only after Python and DejaVu agree; the next optimizer proposal is authorized
before the only state mutation; and independent replay must reconstruct the
same three-level verdict.  The result is a bounded, evidence-backed assurance
argument for this concrete execution, with its trust assumptions and scaling
limits stated explicitly.
